\documentclass{article}
\usepackage{/globals/de}
\usepackage{/globals/math}
\usepackage{/globals/diagrams}
\usepackage{/globals/anki}
\ankiprefix{gatter}
\ankitopic{Gatter}
\ankishow
\begin{document}
\section{Gatter}
Gatter sind die einfachsten logischen Funktionen die aus Transistoren gebaut werden können. Sie können zu Schaltnetzen zusammengesetzt werden.

\subsection{Darstellung}
Im Sinne dieses Moduls werden die deutschen Symbole genutzt. Es gibt für jede logische Operation ein Gatter.
\begin{description}
 \item[AND-Gatter] Mit einem $\&$-Symbol nimmt es zwei Eingaben an und gibt aus ob beide wahr sind.
Das Gatter stellt somit $z = a \land b$ dar.
\ztikz{
 \andg(0,0){A}
 \zlabel(Ainput1){left}{$a$}
 \zlabel(Ainput2){left}{$b$}
 \zlabel(Aoutput){right}{$z$}
}
Es kann auch als Gatter mir mehr als zwei Eingaben gezeichnet werden. Hier ist \zb $z = a \land b \land c$ dar.
\ztikz{
 \tandg(0,0){A}
 \zlabel(Ainput1){left}{$a$}
 \zlabel(Ainput2){left}{$b$}
 \zlabel(Ainput3){left}{$c$}
 \zlabel(Aoutput){right}{$z$}
}
 \item[OR-Gatter] Mit einem $\ge 1$-Symbol nimmt es zwei Eingaben an und gibt aus ob mindestens eine der beiden wahr ist.
Das Gatter stellt somit $z = a \lor b$ dar.
\ztikz{
 \org(0,0){A}
 \zlabel(Ainput1){left}{$a$}
 \zlabel(Ainput2){left}{$b$}
 \zlabel(Aoutput){right}{$z$}
}
Es kann auch als Gatter mir mehr als zwei Eingaben gezeichnet werden. Hier ist \zb $z = a \lor b \lor c$ dar.
\ztikz{
 \torg(0,0){A}
 \zlabel(Ainput1){left}{$a$}
 \zlabel(Ainput2){left}{$b$}
 \zlabel(Ainput3){left}{$c$}
 \zlabel(Aoutput){right}{$z$}
}
 \item[NOT-Gatter] Mit einem $1$-Symbol und einem Invertierpunkt nimmt es eine Eingabe an und gibt dessen Gegenteil aus.
Das Gatter stellt somit $z = \neg a$ dar.
\ztikz{
 \notg(0,0){A}
 \zlabel(Ainput){left}{$a$}
 \zlabel(Aoutput){right}{$z$}
}
 \item[XOR-Gatter] Mit einem $=1$-Symbol nimmt es zwei Eingaben an und gibt aus ob genau eine der beiden, aber nicht beide, wahr sind.
Das Gatter stellt somit $z = a \oplus b$ dar.
\ztikz{
 \xorg(0,0){A}
 \zlabel(Ainput1){left}{$a$}
 \zlabel(Ainput2){left}{$b$}
 \zlabel(Aoutput){right}{$z$}
}
 \item[NAND-Gatter] Mit einem $\&$-Symbol und einem Invertierpunkt nimmt es zwei Eingaben an und gibt aus ob nicht beide wahr sind.
Das Gatter stellt somit $z = \neg (a \land b)$ dar.
\ztikz{
 \nandg(0,0){A}
 \zlabel(Ainput1){left}{$a$}
 \zlabel(Ainput2){left}{$b$}
 \zlabel(Aoutput){right}{$z$}
}
 \item[NOR-Gatter] Mit einem $\ge 1$-Symbol und einem Invertierpunkt nimmt es zwei Eingaben an und gibt aus ob keine der beiden wahr ist.
Das Gatter stellt somit $z = \neg (a \lor b)$ dar.
\ztikz{
 \norg(0,0){A}
 \zlabel(Ainput1){left}{$a$}
 \zlabel(Ainput2){left}{$b$}
 \zlabel(Aoutput){right}{$z$}
}
\end{description}
Es kann jeder Zeit an eine Eingabe oder eine Ausgabe ein unausgefüllter Punkt hinzugefügt werden um darzustellen dass dessen Wert dort invertiert wurde. So können Schaltungen mit vielen NOT-Gatten massiv vereinfacht werden.
\ankinote{and}{and}{\(\&\)}
\ankinote{or}{or}{\(\ge 1\)}
\ankinote{not}{not}{1 p}
\ankinote{xor}{xor}{\(=1\)}
\ankinote{nicht}{nicht Input}{Unausgefüllter Punkt}

\end{document}